\chapter{공기를 통한 대화}
% full version: chapters/05_serocomputer.tex

\pencilfig{5}{\pencilart{5}}{왼쪽은 전화: 한 사람에서 한 사람으로 가는 전선. 오른쪽은 라디오: 하나의 메시지가 주변 모두에게 한꺼번에.}

지금까지 뉴런들은 전화로 이야기했다. 전선에서 전선으로, 주소를 정해서. 이제 첫 방송국을 켜자. 세로토닌을 방출하는 \nt{SERO} 뉴런이다. 뇌 전체에 수십만 개뿐이지만, 이들의 전선은 뇌의 넓은 영역으로 뻗어 나간다.

\section*{네 가지 차이}

\nt{SERO} 뉴런은 \nt{GLUT}와 다르다. 첫째, 신호의 부호는 수신자가 정한다. 세로토닌에는 두 부호의 “자물쇠”가 모두 있다. 둘째, 스스로 일한다. 깨어 있을 때 일정한 배경 입력을 받으면 1초에 몇 번씩 고르게 딸깍한다. 메트로놈이다. 셋째, 느리다. 신호가 밀리초가 아니라 수백 밀리초에 걸쳐 작용한다. 넷째, 물질을 좁은 틈이 아니라 주변 공간으로 내보내는 일이 많다. 물질은 물속 잉크 방울처럼 퍼진다. 수신자는 농도를 느낄 뿐, 분자가 누구에게서 왔는지 모른다.

\section*{있을 법하지만 없는 논리}

억제 “자물쇠”를 단 메트로놈은 그 자체로 “아니다”이다. 세로토닌이 없는 동안 뉴런은 일하고, 세로토닌이 나타나면 입을 다문다. 회로 하나 대신 뉴런 하나이다. 세로토닌 네트워크가 논리적으로 \nt{GLUT} 네트워크보다 강해 보인다.

그런데 걸림돌이 있다. 공간으로 방출된 물질은 섞인다. 두 신호가 구별되려면 물질이 퍼지는 거리보다 더 멀리 떨어진 곳에서 방출되어야 한다. 그 거리는 수십 마이크로미터이다. 그런데 세로토닌 뉴런 하나하나에는 넓은 영역에 흩어진 방출 지점이 엄청나게 많고, 서로 다른 뉴런의 “전선”이 겹친다. 이런 네트워크로는 개별 비트를 보낼 수 없다. 모두가 같은 방송국을 듣는 라디오처럼, 몇 가지 공통의 양만 전할 수 있다.

\pencilfig{123}{%
% two release sites: far apart the clouds stay separate (two signals), close together they merge (one level)
\foreach \x/\y in {0.9/1.55, 4.1/1.55, 1.9/0, 2.9/0}{
  \foreach \r/\o in {0.95/0.07, 0.7/0.09, 0.45/0.12}{\fill[pcViolet, opacity=\o] (\x,\y) circle (\r);}
  \draw[dotpen=pcViolet, fill=pcViolet] (\x,\y) circle (0.07);}
\draw[arr] (1.95,1.55) -- (3.05,1.55); \draw[arr] (3.05,1.55) -- (1.95,1.55);
\node[lbl, anchor=south, align=center] at (2.5,1.6) {퍼지는 거리보다\\멀리};
\node[lbl, anchor=west] at (5.25,1.55) {두 신호};
\node[lbl, anchor=west] at (5.25,0) {하나의 공통 구름};
}{세로토닌은 공간으로 방출되어 구름처럼 퍼진다. 두 방출원은 물질이 퍼지는 거리보다 멀리 떨어져 있을 때만 서로 다른 두 신호를 낸다. 더 가까우면 구름이 합쳐져 하나의 공통 수준이 된다.}

\section*{실제로 하는 일}

공통의 양 하나를 다루는 데 필요한 것은 이 시스템에 다 있다. 세로토닌 농도는 개별 딸깍을 매끄러운 수준으로 평활한다. 가장 단순한 변환기가 스파이크 열을 전압으로 바꾸는 것과 같다. 이 수준은 몇 초 동안 유지되다가 저절로 사라진다. 지워야 하는 비트가 아니라 천천히 사라지는 흔적이다.

세로토닌 뉴런에는 자기 자신에게 달린 “자물쇠”도 있다. 방출된 세로토닌이 제 방출원을 조금 억제한다. 엔지니어는 여기서 음의 되먹임 증폭기, 즉 수준을 유지하는 온도 조절기를 알아볼 것이다. 다만 실험을 보면 이 자기 억제는 특정 대상을 향해 있고 짧은 멈춤만 일으킨다. 그래서 온도 조절기 역할은 아직 사실이라기보다 가설에 가깝다.

\section*{기다림의 값}

네트워크 전체에 대해 느리고 똑같이 유지할 만한 공통의 양은 무엇일까? 좋은 후보는 \emph{인내심}이다. 이런 선택을 떠올려 보자. 지금 받는 작은 보상인가, 몇 초 뒤에 받는 큰 보상인가. 얼마나 기다릴 수 있는지는 미래가 여러분에게 얼마나 빨리 “평가절하”되는지에 달려 있다. 한 가설에 따르면 바로 이 다이얼을 세로토닌이 돌린다. 이를 뒷받침하는 실험도 있다. 세로토닌 뉴런을 인위적으로 켜면 생쥐는 포기하기 전까지 지연된 보상을 더 오래 기다린다.

이 장은 뇌의 모든 방송국이 쓸 세 가지 장치를 소개한다. 메트로놈 뉴런, 실제로는 부피인 “전선”, 그리고 개별 딸깍 대신 쓰는 느린 수준이다.

\fullversion{5}
